% === Thread 1: Instruction Following & SFT ===
\section{Thread 1: Instruction Following \& SFT}
\label{sec:t1}

\begin{quote}
\textit{Core Question: How can we teach a pre-trained language model to understand and follow human instructions?}
\end{quote}

Instruction following is the earliest and most fundamental objective of post-training.
A language model that has only undergone pre-training is essentially a next-token predictor:
given the prompt ``What is the capital of France?'', it may continue with ``What is the capital of Germany?''
rather than answering ``Paris''.
Supervised Fine-Tuning (SFT) trains the model on (instruction, response) data pairs in a supervised manner,
teaching it to understand the intent behind instructions and generate outputs that conform to the expected format.

This chapter traces the complete evolutionary chain of the SFT field from problem formulation to current open problems.
We will see how this field has progressed along two core axes:
\textbf{the data dimension} (from human annotation to synthetic generation, from pursuing quantity to emphasizing quality)
and \textbf{the task dimension} (from single-task to multi-task scaling, from general capabilities to domain-specific specialization).

\evolution{Pre-trained LM can only continue text, not answer questions}{Instruction-tuned LM understands instructional intent}{SFT on (instruction, response) pairs}


% ========================================================
\subsection{Problem Origins: Why Is Pre-Training Not Enough?}
\label{sec:t1-origin}

The objective function of pre-training---whether autoregressive language modeling or masked language modeling---exhibits
a fundamental misalignment with how users actually interact with the model.
This problem can be understood from three perspectives:

\begin{problembox}
\begin{enumerate}[leftmargin=1.5em]
  \item \textbf{Objective function mismatch}:
    Pre-training optimizes $\mathcal{L}_{\text{PT}} = -\sum_t \log P(x_t \mid x_{<t})$,
    i.e., the log-likelihood of predicting every token in the training corpus.
    This objective treats all tokens equally, without distinguishing between instruction comprehension, knowledge recall, and format adherence.
    However, what users need is: given an instruction $x$, the model should generate a response $y$ that is helpful, accurate, and safe.

  \item \textbf{Behavioral mode mismatch}:
    A pre-trained model learns the most common text distribution in web-scale corpora.
    This means it mimics various styles found in the training data---news articles, social media posts, code comments---rather
    than focusing on the conversational mode of an AI assistant.
    In other words, a pre-trained model knows a great deal, but does not know how to express it in the capacity of an assistant.

  \item \textbf{Format mismatch}:
    A pre-trained model lacks awareness of output formatting.
    When a user asks ``List 5 benefits of exercise'', the model may generate a prose paragraph rather than a numbered list;
    when a user asks ``Is X true?'', the model may continue generating more questions rather than providing an answer.
\end{enumerate}
\end{problembox}

\noindent
The core idea of instruction tuning is to perform supervised fine-tuning on a small but high-quality set of (instruction, response) data,
teaching the model to:
(1) extract user intent from the instruction,
(2) invoke the knowledge acquired during pre-training to respond, and
(3) organize the output in the expected format and style.

\begin{solutionbox}
Given an instruction-response dataset $\mathcal{D} = \{(x_i, y_i)\}_{i=1}^N$,
the SFT training objective is:
\[
  \mathcal{L}_{\text{SFT}} = -\sum_{i=1}^{N} \sum_{t=1}^{|y_i|}
    \log P_\theta(y_{i,t} \mid x_i, y_{i,<t})
\]
Note that the loss is computed only over the response tokens $y$ (the instruction tokens $x$ are masked),
ensuring that the model learns ``how to respond to instructions'' rather than ``how to generate instructions''.
\end{solutionbox}


It is worth noting that instruction tuning is not the only pathway for adapting pre-trained LMs.
Prompt tuning \citep{lester2021prompt} and Prefix-Tuning \citep{li2021prefix}
adopt an alternative approach---freezing the model parameters and optimizing only continuous prompt embeddings at the input.
These parameter-efficient methods perform well on specific tasks
but struggle to endow the model with general-purpose instruction following capabilities.
Subsequent experiments in FLAN (Fig.~10) even demonstrated that instruction tuning can \emph{facilitate}
the effectiveness of subsequent prompt tuning, suggesting that the two approaches are complementary rather than competing.

% ========================================================
\subsection{Foundational Methods: Early Explorations and Paradigm Establishment}
\label{sec:t1-foundation}

Between 2021 and 2022, three independent and nearly simultaneous works demonstrated
the effectiveness of instruction tuning from different angles, establishing the foundational paradigm for this field.


% --- FLAN ---
\subsubsection{FLAN: Instruction Tuning for Zero-Shot Generalization}
\label{sec:t1-flan}

\landmark{wei2022flan}
FLAN (Finetuned Language Net) was one of the earliest works to systematically validate
that instruction tuning can improve the zero-shot generalization capabilities of models.

\paragraph{Core design.}
FLAN's experiments were based on LaMDA-PT (a 137B-parameter decoder-only Transformer),
using 62 NLP datasets organized into 10 clusters by task type
(e.g., Natural Language Inference, Sentiment Analysis, Reading Comprehension, etc.).
For each dataset, the authors manually wrote 10 instruction templates to convert traditional NLP samples into
natural language instruction format.
For example, an NLI sample (\texttt{premise}, \texttt{hypothesis}, \texttt{label})
was converted into ``Premise: [P]. Hypothesis: [H]. Does the premise entail the hypothesis?''.

\paragraph{Key experimental findings.}
FLAN's experimental design was remarkably rigorous---when evaluating zero-shot performance on a given cluster,
all tasks from that cluster were \textbf{excluded} from the training data, ensuring that genuine cross-task generalization was being measured:

\begin{enumerate}[leftmargin=1.5em]
  \item \textbf{Zero-shot FLAN surpasses zero-shot GPT-3}:
    On the majority of the 25 evaluation tasks, FLAN's zero-shot performance exceeded that of GPT-3's zero-shot performance,
    and on some tasks even approached or surpassed GPT-3's few-shot performance.

  \item \textbf{Scaling effect of task cluster count}:
    As the number of task clusters included in training increased from 1 to 7,
    the average performance on held-out clusters rose from 49.9\% to 63.5\%,
    exhibiting a nearly linear growth trend.
    This demonstrated that the generalization capability of instruction tuning increases monotonically with task diversity.

  \item \textbf{Threshold effect of model scale}:
    Instruction tuning yielded significant improvements for large models ($\geq$ 100B parameters),
    but could actually cause performance degradation for small models ($\leq$ 8B parameters).
    The authors hypothesized that this is because small models lack sufficient capacity to learn instruction following across multiple tasks simultaneously,
    leading to catastrophic forgetting.

  \item \textbf{Critical role of instructions}:
    Ablation experiments showed that using natural language instruction templates (55.2\%)
    versus a baseline without instructions (37.3\%) resulted in a gap of nearly 18 percentage points,
    demonstrating that the instructions themselves---not merely multi-task training---are the key driver of improved generalization.
\end{enumerate}

\paragraph{Significance and limitations of FLAN.}
FLAN established the basic paradigm of instruction tuning:
by training on instruction-formatted data from multiple tasks, the model can generalize to unseen tasks.
However, FLAN also had notable limitations:
(1) all instruction templates were manually authored, making scaling costly;
(2) task types were limited to traditional NLP benchmarks, lacking open-domain dialogue and creative generation scenarios representative of real user needs;
(3) validation was conducted only at the 137B scale, leaving applicability to smaller models uncertain.


% --- T0 ---
\subsubsection{T0: Prompted Training for Zero-Shot Generalization}
\label{sec:t1-t0}

\landmark{sanh2022multitask}
T0 was published nearly simultaneously with FLAN, independently validating similar core ideas
while differing in experimental design and model choice.

T0 was based on the encoder-decoder architecture of T5 (11B parameters),
using the Public Pool of Prompts (P3) dataset---comprising
2,073 manually written prompt templates across 62 datasets.
On 11 held-out tasks, T0's zero-shot performance
matched or surpassed GPT-3's few-shot (16-shot) performance on 7 tasks,
despite T0 having only $1/16$ the parameters of GPT-3.

The complementarity between T0 and FLAN is reflected in:
(1) T0 demonstrated that encoder-decoder architectures also benefit from instruction tuning, not just decoder-only models;
(2) the large-scale prompt template library of P3 provided a valuable data resource for subsequent research;
(3) T0's prompt engineering practices accumulated empirical knowledge about how template design affects performance.

Concurrent with FLAN, \citet{mishra2021crossfit} proposed a crowdsourcing-based approach
for collecting cross-task instruction data,
and \citet{wang2022supernaturalinstructions} extended this idea to
Super-Natural\-Instructions---a large-scale benchmark containing 1,616 tasks
that provided important data infrastructure for subsequent multi-task instruction tuning.


% --- InstructGPT ---
\subsubsection{InstructGPT: Establishing the Three-Stage Training Paradigm}
\label{sec:t1-instructgpt}

\landmark{ouyang2022instructgpt}
While FLAN and T0 demonstrated the effectiveness of instruction tuning on academic benchmarks,
InstructGPT extended this approach to real-world user scenarios
and established the highly influential ``SFT $\rightarrow$ Reward Model $\rightarrow$ PPO'' three-stage training paradigm.

\paragraph{Technical details of the SFT stage.}
InstructGPT's SFT stage used approximately 13,000 demonstration data entries written by human annotators.
These data came from two sources:
(1) real user prompts collected from the OpenAI API (filtered and anonymized), and
(2) prompts proactively composed by annotators based on the task distribution.
Training configuration: 16 epochs, cosine learning rate decay, dropout on residual connections ($p=0.2$).
The authors found that multi-epoch training on a small dataset was effective---despite the training data being far smaller than the pre-training corpus,
SFT could still significantly alter the model's behavioral patterns.

\paragraph{Core findings.}
InstructGPT's most striking conclusion was:
\textbf{a 1.3B-parameter SFT+RLHF model surpassed the 175B-parameter original GPT-3 in human evaluation with an 85$\pm$3\% win rate}.
This means that post-training can achieve a qualitative leap without increasing model size---despite
a parameter count difference of over 100$\times$, the smaller model after alignment was preferred by users.
Furthermore, InstructGPT's performance on truthfulness was approximately twice that of GPT-3 (TruthfulQA),
with the hallucination rate dropping from 41\% to 21\%.

\paragraph{Alignment Tax.}
InstructGPT was the first to systematically identify the \textbf{alignment tax}---the
phenomenon where post-training, while improving instruction following and safety,
may cause performance degradation on traditional NLP benchmarks.
The authors mitigated this issue by incorporating pre-training data into the PPO objective through mixed training (PPO-ptx):
\[
  \text{objective}(\theta) =
    \mathbb{E}_{(x,y) \sim \pi_{\text{RL}}} \big[ r_\theta(x,y) - \beta \log \frac{\pi_{\text{RL}}(y|x)}{\pi_{\text{SFT}}(y|x)} \big]
    + \gamma \, \mathbb{E}_{x \sim \mathcal{D}_{\text{pretrain}}} \big[ \log \pi_{\text{RL}}(x) \big]
\]
where the second term alleviates the alignment tax by maintaining the model's performance on the pre-training distribution.
\crossref{2}{sec:t2}

\paragraph{Impact of InstructGPT on subsequent research.}
The impact of InstructGPT extends far beyond its technical contributions.
It established two paradigms that persist to this day:
(1) SFT as a necessary precursor stage of the RLHF/DPO pipeline
(virtually all subsequent preference optimization work starts from an SFT checkpoint);
(2) human evaluation as the gold standard for evaluating instruction-following models
(rather than relying solely on automatic metrics).


% ========================================================
\subsection{Limitations and Branching: Data Challenges and Divergent Solution Paths}
\label{sec:t1-branching}

FLAN, T0, and InstructGPT established the effectiveness of instruction tuning,
but also exposed a core bottleneck: \textbf{the cost of acquiring high-quality instruction data is prohibitively high}.
InstructGPT's 13,000 demonstration data entries were completed by 40 professional annotators over several months,
a cost that is unsustainable for the academic community and the open-source ecosystem.

\begin{problembox}
How can large-scale, high-quality instruction tuning data be obtained at low cost?
\end{problembox}

This problem gave rise to two nearly opposite solution paths:
\textbf{Branch A} pursues data quantity expansion (reducing costs through synthetic generation),
while \textbf{Branch B} pursues data quality refinement (achieving more with less through careful curation).
The tension between these two approaches constitutes the most central debate in the SFT field.

\evolution{High cost of human-annotated data}{Branch A: Synthetic data / Branch B: Data curation}{Lowering the barrier to obtaining instruction data}


% --- Branch A: Synthetic Data ---
\subsubsection{Branch A: Synthetic Data Generation---Winning Through Quantity}
\label{sec:t1-synthetic}

The idea of using language models themselves to generate training data can be traced back to
\citet{schick2021generating}, who demonstrated that GPT-2-level models
could generate effective labeled data for SuperGLUE tasks,
and that models trained on generated data approached the performance of those trained on human-annotated data for some tasks.
This result foreshadowed the possibility that ``an LM can serve as both student and teacher''.

\paragraph{Self-Instruct: Bootstrapping instruction data from the LM itself.}
\landmark{wang2023selfinstruct}
Self-Instruct extended this approach from task-specific labeled data to the generation of \textbf{general-purpose instruction data},
becoming the most influential work in the synthetic instruction data domain.
It pioneered the bootstrap paradigm of leveraging a model's own capabilities to generate training data.

\textbf{Pipeline design.}
The Self-Instruct pipeline consists of four stages:
\begin{enumerate}[leftmargin=1.5em]
  \item \textbf{Instruction generation}:
    Starting from 175 manually written seed tasks, new instructions are generated by
    vanilla GPT-3 (text-davinci-001) through in-context learning.
    Each round samples 6 human seeds + 2 model-generated instructions from the task pool as few-shot examples.

  \item \textbf{Classification identification}:
    Each instruction is classified as either a classification task or a non-classification task,
    as the two require different instance generation strategies.

  \item \textbf{Instance generation}:
    For non-classification tasks, an input-first approach is used (generating input before output);
    for classification tasks, an output-first approach is used (determining the label before generating input),
    which effectively prevents label distribution skew.

  \item \textbf{Filtering}:
    A ROUGE-L $< 0.7$ threshold is applied to filter out newly generated samples that are overly similar to existing data,
    ensuring data diversity.
\end{enumerate}

\textbf{Results.}
Starting from 175 seed tasks, Self-Instruct generated 52,445 instructions and 82,439 instances.
Quality evaluation showed: 92\% of generated instructions were valid, and 54\% of samples were valid across all fields.
On the SuperNI benchmark, GPT-3 fine-tuned with Self-Instruct data
outperformed vanilla GPT-3 by 33.1\% ROUGE-L,
and in human evaluation fell only about 5\% behind InstructGPT$_{\text{001}}$.
The total API cost for the entire data generation process was approximately \$600.

\textbf{Significance.}
The significance of Self-Instruct lies in demonstrating that \textbf{an LM can serve as a source of instruction data itself},
an idea that directly inspired a series of subsequent works.

\paragraph{Alpaca: Industrializing Self-Instruct.}
Stanford Alpaca \citep{taori2023alpaca} applied the Self-Instruct approach to LLaMA 7B,
using GPT-3.5 (text-davinci-003) to generate 52K instruction data entries,
with a total fine-tuning cost of only approximately \$600 (API) + \$100 (training).
Alpaca demonstrated performance comparable to GPT-3.5 in human evaluation,
marking a milestone for open-source instruction-following models.
However, Alpaca also exposed the limitations of synthetic data:
the generated data contained significant hallucination and style mimicry (imitating the teacher model's linguistic style
rather than truly learning the underlying capabilities), sparking discussion about the ceiling of synthetic data.
Concurrently, Vicuna \citep{chiang2023vicuna} offered a contrasting path:
it used approximately 70K \textbf{real user conversations with ChatGPT} collected from the ShareGPT platform
to fine-tune LLaMA 13B, achieving approximately 90\% of ChatGPT's quality in GPT-4 evaluation.
Vicuna's success demonstrated that real conversation data holds advantages
that synthetic data can hardly replicate in capturing user interaction patterns---the way users phrase questions, contextual dependencies in multi-turn dialogue,
and implicit preferences are all embedded in real interactions.
The parallel development of Alpaca (synthetic data) and Vicuna (real conversation data)
laid the empirical foundation for subsequent discussions about data source selection.

\paragraph{WizardLM: Evol-Instruct for progressive complexity.}
\citet{xu2024wizardlm} proposed the Evol-Instruct method---rather
than generating instructions from scratch, it iteratively evolves existing simple instructions.
In each evolution round, a prompt directs the LLM to perform one of the following operations on an instruction:
adding constraints, increasing reasoning steps, concretizing scenarios, or increasing input complexity.
After multiple rounds of evolution, a simple ``Write a poem about spring'' might become
``Write a Shakespearean sonnet about spring, incorporating at least three
metaphors related to scientific phenomena, and ending with a volta
that challenges the reader's expectations''.
This method achieves controlled data augmentation along the instruction complexity dimension.

\paragraph{Other synthetic data works.}
\citet{honovich2022unnatural}'s Unnatural Instructions adopted an approach similar to Self-Instruct
but used fewer seeds (15), validating the robustness of bootstrapping.
\citet{peng2023gpt4instruct} used GPT-4 to generate higher-quality instruction data,
demonstrating the decisive influence of teacher model capability on synthetic data quality.
Orca \citep{mukherjee2023orca} proposed a progressive learning strategy---learning
from GPT-4's detailed reasoning traces,
not only imitating the final answer but also the reasoning process (the system message instructs GPT-4 to ``explain step by step''),
an idea that foreshadowed the development of reasoning distillation in Thread 4
\crossref{4}{sec:t4}.
Baize \citep{xu2023baize} explored the use of ChatGPT self-dialogue to generate multi-turn conversation data.


% --- Branch B: Data Quality ---
\subsubsection{Branch B: Data Quality Curation---Winning Through Quality}
\label{sec:t1-quality}

\paragraph{LIMA: The Superficial Alignment Hypothesis.}
\landmark{zhou2023lima}
LIMA proposed one of the most controversial hypotheses in the SFT field---the
\textbf{Superficial Alignment Hypothesis}:
\textit{a model's knowledge and capabilities are acquired almost entirely during pre-training,
and the role of alignment is merely to teach the model the correct output format and interaction style}.

\textbf{Experimental design.}
LIMA fine-tuned LLaMA 65B using only 1,000 carefully curated training samples.
The diversity of data sources was meticulously designed:
750 samples came from community Q\&A platforms (highest-quality responses from Stack Exchange,
wikiHow articles, highly upvoted Reddit posts),
and 250 were manually written by the authors (ensuring coverage of specific conversational abilities and safety behaviors).
The training configuration appeared simple but was carefully tuned:
15 epochs, AdamW optimizer, residual dropout increasing linearly from $p_d = 0.0$ at the bottom layers to $p_d = 0.3$ at the top layers.
This \textbf{layer-wise increasing dropout} is a noteworthy detail---it
protects the representational stability of lower layers (closer to pre-training knowledge)
while allowing greater learning flexibility in upper layers (closer to output formatting).

\textbf{Core results.}
In human evaluation, LIMA's performance was striking:
\begin{itemize}[leftmargin=1.5em]
  \item 50\% of outputs were rated ``excellent'', 88\% ``pass'' or better
  \item Against Alpaca 65B: LIMA was preferred in 53\% of cases
  \item Against DaVinci003: LIMA was preferred or tied in 44\% of cases
  \item Against GPT-4: LIMA was comparable or superior in 43\% of cases
\end{itemize}

\textbf{Ablation experiments} (conducted on 7B models) provided more granular insights:
\begin{itemize}[leftmargin=1.5em]
  \item \textbf{Diversity $>$ homogeneity}:
    Diverse Stack Exchange data (rating 3.83/6) outperformed homogeneous wikiHow data (3.49/6),
    demonstrating that source diversity matters more than data volume.
  \item \textbf{Quality filtering is effective}:
    Manually filtered data (3.83/6) significantly outperformed unfiltered data (3.33/6).
  \item \textbf{Quantity scaling has a ceiling}:
    After increasing from 2K to 4K samples, performance gains plateaued.
  \item \textbf{A small amount of dialogue data has outsized impact}:
    Adding just 30 multi-turn dialogue samples raised the ``excellent'' rating from 45.2\% to 76.1\%.
\end{itemize}

\textbf{Controversies surrounding LIMA.}
LIMA's Superficial Alignment Hypothesis sparked ongoing debate.
Proponents argue that it reveals the essence of SFT---not teaching the model new knowledge, but activating the correct expression of existing capabilities.
Critics counter that:
(1) it remains unclear whether the conclusions from 1,000 samples generalize to more complex capabilities (e.g., multi-step reasoning, code generation);
(2) LIMA used a 65B-parameter LLaMA---for smaller models, SFT may genuinely need to impart new knowledge rather than merely adjust formatting;
(3) LIMA's evaluation was primarily based on single-turn open-domain Q\&A, with insufficient coverage of specialized domains and long-horizon tasks.

\paragraph{Data selection methodologies.}
LIMA's success inspired a series of data selection methods.
\citet{li2023ifd} proposed the Instruction-Following Difficulty (IFD) metric---measuring
the informativeness of each data sample by comparing the difference in the model's output probability with and without the instruction,
then selecting samples with the highest IFD for training.
DEITA \citep{liu2024deita} proposed an automated data scoring framework based on three dimensions---complexity, quality, and diversity---using ChatGPT to score and rank data,
then selecting the optimal subset through embedding-based diversity filtering.
\citet{chen2023maybeonly} found that using only 0.5\% of the Alpaca data could achieve over 90\% of the performance of full-data training.
Instruction Mining \citep{cao2023instructionmining} trained a quality predictor
to automate the data filtering process.

\paragraph{Open-source datasets and community contributions.}
OpenAssistant Conversations \citep{kopf2023openassistant} provided
a multilingual dialogue dataset contributed by 13,500 volunteers (161,443 messages in 35 languages),
demonstrating the potential of crowdsourcing for instruction data collection.
Dolly \citep{conover2023dolly}, contributed by 5,000 Databricks employees during work hours with approximately 15,000 entries,
was one of the earliest enterprise-level open-source instruction datasets.


% --- Multi-task Scaling ---
\subsubsection{Branch C: Multi-task Scaling---Expansion Along the Task Quantity Dimension}
\label{sec:t1-multitask}

\paragraph{Flan-v2: Pushing instruction tuning to the extreme.}
\landmark{chung2022flan}
Flan-v2 is the large-scale follow-up to FLAN,
pushing the task scale of instruction tuning to unprecedented levels.

\textbf{Data scale.}
Flan-v2 integrated four major task collections:
Muffin (80 tasks), T0-SF (193 tasks), NIV2 (1,554 tasks), and CoT (9 reasoning tasks),
totaling 1,836 tasks (473 unique datasets, 146 task categories).
This represents approximately a 30$\times$ expansion over the original FLAN's 62 tasks.

\textbf{Core findings.}
\begin{enumerate}[leftmargin=1.5em]
  \item \textbf{Cross-architecture effectiveness}:
    Flan-v2 validated its effectiveness on T5 (encoder-decoder, 80M--11B),
    PaLM (decoder-only, 8B--540B), as well as cont-PaLM and U-PaLM,
    demonstrating that the benefits of instruction tuning are architecture-agnostic.

  \item \textbf{Diminishing marginal returns of task scaling}:
    Performance continued to improve with increasing task count, but the majority of gains were concentrated in the first 282 tasks,
    with diminishing marginal contributions from subsequent tasks.
    This suggests an important practical insight:
    achieving ``good enough'' performance does not require exhausting every possible task.

  \item \textbf{Extreme efficiency advantages}:
    Taking PaLM 540B as an example, Flan fine-tuning used only 512 v4 TPU chips for 37 hours,
    amounting to a mere \textbf{0.2\%} of the pre-training compute.
    Yet the effects were substantial: Flan-PaLM 540B achieved 75.2\% on MMLU (5-shot, with CoT + self-consistency),
    with a normalized average improvement of 9.4--15.5\%.
    Even more remarkably, \textbf{Flan-T5-XL 3B achieved 52.4\% on MMLU, surpassing GPT-3 175B's 43.9\%}---a
    model with $1/58$ the parameters outperformed a $58\times$ larger model after instruction tuning.

  \item \textbf{Critical role of CoT fine-tuning}:
    An underappreciated yet highly important finding in Flan-v2 concerns Chain-of-Thought fine-tuning.
    Instruction tuning without CoT data actually \textbf{harmed} the model's reasoning capabilities;
    however, jointly training with both CoT and non-CoT data
    not only preserved reasoning capabilities but also improved performance on both types of tasks.
    Furthermore, CoT fine-tuning unlocked zero-shot CoT capability---merely
    appending ``Let's think step by step'' to the end of a prompt
    could elicit reasoning behavior from the model, without requiring few-shot reasoning examples.
    This finding foreshadowed the development of reasoning post-training in Thread 4
    \crossref{4}{sec:t4}.
\end{enumerate}

The Flan Collection \citep{longpre2023flancollection} subsequently open-sourced the complete data mixing recipe and templates of Flan-v2,
providing the community with a reproducible foundation.
OPT-IML \citep{iyer2022optiml} replicated similar multi-task instruction tuning experiments on Meta's OPT model.
BLOOMZ/mT0 \citep{muennighoff2022bloomz} extended multi-task instruction tuning
to multilingual scenarios (crosslingual instruction tuning with the xP3 dataset),
demonstrating the effectiveness of this paradigm for non-English languages.


% --- Domain-specific SFT ---
\subsubsection{Branch D: Domain-Specific SFT---Precise Adaptation for Specialized Domains}
\label{sec:t1-domain}

As general-purpose instruction tuning matured, researchers began applying SFT techniques to
vertical domains requiring specialized knowledge, sparking a wave of domain-specific SFT research.

\paragraph{Code: Data quality can break scaling laws.}
\landmark{gunasekar2023phi1}
phi-1 is the most disruptive work in the domain of code SFT.
This research from Microsoft Research used a model of only 1.3B parameters
to achieve 50.6\% pass@1 on HumanEval and 55.5\% pass@1 on MBPP---at
the time surpassing many models an order of magnitude larger.

\textbf{The core innovation lies in data engineering.}
phi-1's training data comprises three carefully designed components:
\begin{itemize}[leftmargin=1.5em]
  \item \textbf{Filtered code-language corpus} (approximately 6B tokens):
    A random forest classifier was used to filter ``textbook quality'' code snippets
    from The Stack and StackOverflow.
    The classifier was trained on a small set of GPT-4-labeled samples (annotating the educational value of each code snippet),
    then applied to filter the entire corpus at scale.

  \item \textbf{Synthetic textbooks} ($< 1$B tokens):
    GPT-3.5 was used to generate textbook-style Python teaching content,
    including concept explanations, code examples, and progressive exercises.

  \item \textbf{CodeExercises} ($< 180$M tokens):
    GPT-3.5 was used to generate function-level programming exercises (function docstring + solution),
    serving as training data for the fine-tuning stage.
\end{itemize}

\textbf{Key insights.}
phi-1's success challenged the conventional scaling law mindset---it
demonstrated that \textbf{data quality can, to a certain extent, substitute for model scale and data volume}.
A 10$\times$ smaller model trained on 100$\times$ less data
could match or even surpass the performance of larger models.
More interestingly, phi-1 exhibited some \textbf{emergent capabilities} after fine-tuning---including
the ability to use external libraries such as Pygame and Tkinter as well as chat-style interaction modes---capabilities
that were not present in the fine-tuning data, suggesting that high-quality data may help the model
better organize and generalize the knowledge acquired during pre-training.
This finding resonates intriguingly with LIMA's Superficial Alignment Hypothesis.

WizardCoder \citep{luo2023wizardcoder} applied WizardLM's Evol-Instruct method to the code domain,
generating high-difficulty training data through progressive complexification of programming instructions.
Magicoder \citep{wei2023magicoder} proposed OSS-Instruct---starting
from real open-source code snippets, it has the LLM generate corresponding instructions for these code segments,
ensuring that the code style and complexity of training data are consistent with real-world projects.

\paragraph{Math: Reasoning-enhanced SFT.}
Mathematical reasoning is another active direction in domain-specific SFT.
WizardMath \citep{luo2023wizardmath} extended Evol-Instruct to the mathematics domain,
combining Reinforced Evol-Instruct to further improve data quality.
MAmmoTH \citep{yue2023mammoth} constructed a hybrid instruction dataset mixing CoT and Program-of-Thought (PoT)---using
Python code as reasoning steps for problems amenable to symbolic computation,
and natural language CoT for problems requiring semantic understanding.
MetaMath \citep{yu2023metamath} proposed question bootstrapping for mathematical problems---augmenting
data diversity by rephrasing problem statements (synonym substitution, condition inversion, self-questioning decomposition)
while maintaining mathematical correctness.
RFT \citep{yuan2023rft} (Rejection sampling Fine-Tuning)
investigated the sampling-then-filtering strategy for mathematical reasoning:
for each math problem, the model generates multiple candidate answers, retaining only samples with correct final answers for SFT,
and analyzed the scaling relationship between sampling quantity and final performance.

\paragraph{Multilingual \& Other Domains.}
Aya Model \citep{ustun2024aya} is a significant milestone in multilingual instruction tuning,
covering 101 languages and demonstrating both the feasibility and challenges of instruction tuning in extremely multilingual settings.
Gorilla \citep{patil2023gorilla} focused on API calling capabilities,
performing SFT on API documentation and calling examples to enable models to accurately generate complex API calling code
\crossref{6}{sec:t6}.
Platypus \citep{lee2023platypus} conducted small-scale curated SFT for the STEM domain.
LongForm \citep{koksal2023longform} focused on instruction tuning for long-form text generation tasks.


% ========================================================
\subsection{Convergence and Frontiers: Systematic Synthesis and Methodological Fusion}
\label{sec:t1-convergence}

In the second half of 2023, the SFT field transitioned from a period of diverse exploration to a phase of systematic synthesis.
Multiple works began integrating the experience from the aforementioned branches, seeking to answer:
\textbf{How can data engineering, multi-task training, and domain specialization
be fused into a unified set of best practices?}

\evolution{Branch A (synthetic) / Branch B (curation) / Branch C (scaling) developing independently}%
  {Systematic synthesis: best practices fusing multiple paths}{Experiment-driven recipe research}


% --- Tulu ---
\subsubsection{Tulu: Systematic Exploration of Open-Source SFT Recipes}
\label{sec:t1-tulu}

Tulu \citep{wang2023tulu} was the first work to systematically compare different instruction data
sources and training strategies in an open-source setting.
Through controlled experiments on LLaMA, Tulu answered multiple practical questions:
Which combinations of data sources are most effective? How does the data mixing ratio affect performance?
What are the relative advantages of human-written and synthetic data across different tasks?

Tulu 2 \citep{ivison2023tulu2} further expanded this research,
extending the experimental scope from SFT to the full post-training pipeline (SFT + DPO)
and validating optimal recipes on the LLaMA 2 series of models.
Tulu 2's core contribution was providing a reproducible training recipe---including
data mixing strategies, hyperparameter configurations, and evaluation protocols---enabling
the open-source community to iteratively optimize on this foundation.

\subsubsection{Zephyr \& OpenChat: Knowledge Distillation from Strong to Weak Models}
\label{sec:t1-distillation}

Zephyr \citep{tunstall2023zephyr} proposed the Direct Distillation approach:
first using GPT-4 as a teacher model to generate high-quality instruction data for SFT,
then using GPT-4's preference annotations for DPO training.
The entire process required no human annotation whatsoever, achieving end-to-end
knowledge distillation from a proprietary model to an open-source model.
Zephyr-7B outperformed most 70B-scale open-source models on MT-Bench at the time,
demonstrating the extreme effectiveness of distillation + alignment for small models
\crossref{7}{sec:t7}.

OpenChat \citep{wang2023openchat} addressed another practical issue:
when training data comes from mixed-quality sources (e.g., some from GPT-4, some from GPT-3.5),
how can such data be effectively utilized?
OpenChat proposed Conditioned RLFT (C-RLFT)---a
class-conditioned training strategy that enables the model to distinguish and leverage data of varying quality,
rather than simply discarding low-quality data as noise.


\subsubsection{Scaling Laws and Theoretical Analysis}
\label{sec:t1-scaling}

\citet{muennighoff2023scaling} systematically investigated
the scaling behavior of instruction-tuned LMs in data-constrained scenarios.
When available instruction data is limited (an extremely common situation in practice),
should one train for more epochs on the limited data (repeat), or augment the limited data and train for one epoch?
Their experiments demonstrated that data repetition training exhibits clear diminishing returns---performance
growth slows significantly beyond 4 epochs,
providing theoretical support for the epoch count choices made by small-data approaches such as LIMA.

The Instruction Tuning Survey \citep{zhang2023instruction} provided
a systematic review of the entire instruction tuning field, tracing the development from FLAN to LIMA
and summarizing best practices in data construction, training strategies, and evaluation methods.


\subsubsection{2025: SFT as the Cornerstone of Reasoning Model Training Pipelines}
\label{sec:t1-2025}

The reasoning model wave of 2025 (\crossref{4}{sec:t4-thinking-models})
endowed SFT with a new role:
SFT is no longer merely the final step that ``teaches the model to follow instructions'',
but has become the \textbf{infrastructure layer} in multi-stage post-training pipelines.

\paragraph{Qwen3's four-stage pipeline.}
Qwen3 \cite{yang2025qwen3}'s four-stage post-training pipeline (detailed in \crossref{4}{sec:t4-thinking-models})
clearly illustrates the multiple roles of SFT in modern training paradigms:
SFT appears simultaneously in Stage 1 (long CoT cold start, reasoning format initialization)
and Stage 3 (thinking mode fusion, integration of reasoning and non-reasoning modes),
marking the evolution of SFT from a ``one-time training step'' to a ``general-purpose alignment tool deployable multiple times within a pipeline''.

\paragraph{Llama 4's natively multimodal SFT.}
Meta's Llama 4 \cite{meta2025llama4} extended the scope of SFT to natively multimodal scenarios,
performing joint instruction tuning on images, text, and code within a unified framework,
further broadening the applicability of SFT methodology
(a systematic discussion of multimodal post-training can be found in \crossref{5}{sec:t5}).

\subsubsection{2026: Token-Level SFT and SFT-RL Synergy}
\label{sec:t1-2026}

In early 2026, the SFT field exhibited two prominent trends:
\textbf{training granularity descending from the sample level to the token level},
and \textbf{SFT being repositioned as the ``optimal initializer'' for RL training}.

\paragraph{Token-Level SFT: From uniform loss to selective training.}
Traditional SFT applies uniform cross-entropy loss to all tokens in a sequence,
but in 2026, multiple independent studies simultaneously challenged this assumption.
Shen et al.\ \cite{shen2026tokenpriority} formalized this emerging paradigm as
\textbf{Token Priority},
pointing out that fine-grained autoregressive generation is constrained by coarse uniform supervision signals,
and categorized existing work into two types:
\emph{Positive Priority} (filtering noisy tokens, retaining high-information tokens)
and \emph{Signed Priority} (applying negative weights to harmful patterns).

ProFit \cite{liu2026profit} provided a concrete instance:
by analyzing the probability distribution of token generation,
it found that \textbf{high-probability tokens carry the core logical framework, while low-probability tokens are mostly substitutable expressions}.
Based on this observation, ProFit selectively retains high-probability tokens for training,
suppressing surface-level overfitting to low-probability, non-core tokens.
TOSS \cite{li2026toss} approached token-level selection from a \textbf{safety perspective}:
by measuring the per-token loss difference between a safety-degraded model and a utility-oriented model,
it quantifies the safety risk of each token,
achieving a fine-grained balance between maintaining task capability and preserving safety
(\crossref{3}{sec:t3}).

\evolution{Sample-level data selection (LIMA, DEITA)}{Token-level selective training (ProFit, TOSS)}{Uniform loss wastes gradient signal on non-critical tokens}

\paragraph{Deepening of data selection methodologies.}
At the sample level, progress in 2026 was equally notable.
NAIT \cite{chen2026nait} proposed a \textbf{neuron activation pattern-based} data selection framework:
by analyzing the \emph{neuron activation similarity} between instruction tuning datasets and target capabilities,
it consistently outperformed methods relying on external strong models using only 10\% of the Alpaca-GPT4 dataset.
NAIT also revealed an intriguing finding:
logical reasoning and programming data exhibit strong generalizability at the neuron level,
enabling cross-capability transfer.
Skill-Aware Data Selection \cite{zhang2026skillaware}
targets reasoning distillation scenarios,
proposing a \textbf{skill-coverage-oriented} data selection strategy:
prioritizing samples corresponding to the target model's weakest skills,
selecting only 1K samples from a 100K corpus to surpass random SFT baselines.

In synthetic data generation,
DS$^2$-Instruct \cite{xu2026ds2instruct} introduced Bloom's taxonomy of cognitive levels
to construct domain-specific instruction datasets:
by pairing domain keywords with different cognitive levels,
it achieved significant zero-shot performance improvements across 7 specialized domains.
Terminal-Task-Gen \cite{pi2026terminaltaskgen} demonstrated
a systematic data engineering pipeline for terminal/CLI tasks,
combining seed-based and skill-based task construction,
data filtering, and curriculum learning strategies,
with the resulting Nemotron-Terminal 32B achieving 27.4\% on Terminal-Bench 2.0.

\paragraph{PEAR: The goal of SFT should be to prepare for RL.}
\landmark{pear2026sftrl}
Perhaps the most paradigm-shifting finding of 2026 comes from PEAR:
\textbf{after identical RL training, a stronger SFT checkpoint may end up significantly weaker than a weaker SFT checkpoint}.
The root cause lies in the \emph{distribution mismatch} between offline SFT data and online RL rollouts---overfitting
during the SFT stage to the teacher data distribution
restricts the exploration space during the RL stage.

PEAR mitigates this distribution mismatch by introducing \textbf{importance sampling}
into the SFT loss (reweighting at the token, block, and sequence levels).
On AIME2025, models initialized with PEAR achieved
pass@8 improvements of up to 14.6\% after RL training.
This result fundamentally redefines the optimization objective of SFT:
\textbf{good SFT is not about maximizing performance on SFT tasks,
but about making the model the best starting point for the RL stage}.

\begin{openproblembox}
\textbf{Open Problem 6: The optimal granularity of SFT.}
The success of token-level selective training indicates that
uniform sequence-level loss may not be optimal.
But what is the optimal training granularity?
Token-level? Span-level? Skill-level?
How do different granularity choices affect downstream RL training?
PEAR's findings suggest that the evaluation criterion for SFT itself needs to be revisited---it
should not be measured by performance on SFT tasks themselves,
but rather optimized for ``RL-readiness'' as the objective.
\end{openproblembox}

% ========================================================
\subsection{Open Problems}
\label{sec:t1-open}

Despite SFT being the most mature technique in post-training, several core problems remain unresolved:

\begin{openproblembox}
\paragraph{Open Problem 1: Where is the ceiling of synthetic data?}
Self-Instruct and Alpaca launched the wave of synthetic instruction data,
but multiple studies have observed that models trained on synthetic data exhibit \textbf{style mimicry without substance}---the
model learns the teacher model's output style, but is brittle when confronting out-of-distribution problems requiring deep understanding.
phi-1 partially mitigated this issue through ``textbook quality'' data filtering,
but a more fundamental question remains: \textbf{does synthetic data have a theoretical capability ceiling?}
That is, can a student model surpass the capability boundary of the teacher model through synthetic data?
This relates to the classic question in knowledge distillation, but in the context of instruction tuning
introduces additional complexity (the transfer of instruction following ability vs. knowledge itself).
\end{openproblembox}

\begin{openproblembox}
\paragraph{Open Problem 2: The optimal interplay between SFT and preference optimization.}
The ``SFT $\rightarrow$ RLHF/DPO'' pipeline established by InstructGPT has become the de facto standard,
but to what extent should the SFT stage be trained?
Excessive SFT (overfitting to instruction data) may limit the exploration space for subsequent RLHF/DPO;
insufficient SFT may lead to instability during the RLHF stage.
Zephyr and Tulu 2's experiments provide some empirical guidance,
but a systematic understanding of the relationship between SFT intensity (training steps, data volume) and downstream alignment effectiveness remains lacking
\crossref{2}{sec:t2}.
\end{openproblembox}

\begin{openproblembox}
\paragraph{Open Problem 3: The capability boundary of instruction following in small models.}
FLAN observed that instruction tuning can cause performance degradation for small models ($\leq$ 8B),
and LIMA's Superficial Alignment Hypothesis implies that SFT effectiveness is highly dependent on pre-training quality.
Then, for models at the 1B--3B scale, where does the capability boundary of instruction following lie?
phi-1 demonstrated impressive code capabilities at the 1.3B scale,
but does this generalize to broader instruction following scenarios?
Understanding the upper efficiency bound of SFT for small models has significant practical implications for edge computing and mobile device scenarios
\crossref{7}{sec:t7}.
\end{openproblembox}

\begin{openproblembox}
\paragraph{Open Problem 4: A systematic solution to catastrophic forgetting.}
SFT inevitably overwrites some knowledge acquired during pre-training,
as manifested concretely in InstructGPT's alignment tax and Flan-v2's reasoning capability degradation due to missing CoT data.
While ad hoc solutions such as PPO-ptx and CoT mixed training are effective in specific scenarios,
a unified continual learning framework with theoretical guarantees is still lacking.
How to precisely control the balance between ``learning new behaviors'' and ``preserving existing knowledge'' during SFT
remains a fundamental cross-thread challenge in the post-training field.
\end{openproblembox}

\begin{openproblembox}
\paragraph{Open Problem 5: Fragmentation of evaluation standards.}
The evaluation of instruction-tuned models is currently highly fragmented:
academia relies on benchmarks (MMLU, SuperNI, HumanEval),
while industry relies on human evaluation and LLM-as-judge approaches (e.g., MT-Bench, AlpacaEval).
The correlations among these evaluation dimensions are weak and unstable---a
model that excels on MMLU may perform mediocrely in human evaluation, and vice versa.
Establishing an evaluation framework that is both scalable and accurately reflects actual user experience
remains one of the most pressing infrastructure needs in this field.
\end{openproblembox}


\subsection{Chapter Summary}
\label{sec:t1-summary}

\begin{figure}[!ht]
\centering
\begin{tikzpicture}[
  node distance=0.6cm and 1.2cm,
  every node/.style={font=\small},
  milestone/.style={rectangle, rounded corners, draw=thread1, fill=thread1!10,
    text width=3.8cm, minimum height=0.8cm, align=center, font=\small\bfseries},
  branch/.style={rectangle, rounded corners, draw=gray!60, fill=gray!5,
    text width=3.2cm, minimum height=0.7cm, align=center, font=\small},
  arrow/.style={-{Stealth[length=2.5mm]}, thick, thread1!70},
  brancharrow/.style={-{Stealth[length=2mm]}, gray!60, thick},
]
% Timeline nodes
\node[milestone] (flan) {FLAN / T0 (2021)\\Multi-task Instruction Tuning};
\node[milestone, below=1.2cm of flan] (igpt) {InstructGPT (2022)\\SFT$\rightarrow$RM$\rightarrow$PPO Pipeline};
\node[milestone, below=1.2cm of igpt] (si) {Self-Instruct (2022)\\Bootstrap from LM};

% Branching
\node[branch, below left=1.2cm and 0.5cm of si] (brA) {Branch A: Synthetic Data\\Alpaca, WizardLM, Orca};
\node[branch, below right=1.2cm and 0.5cm of si] (brB) {Branch B: Data Curation\\LIMA (2023)};

% Convergence
\node[milestone, below=2.8cm of si] (flanv2) {Flan-v2 (2022)\\1836 Tasks Scaling};
\node[milestone, below=1.2cm of flanv2] (phi1) {phi-1 (2023)\\Textbook-Quality Data};
\node[branch, below=1.2cm of phi1] (conv) {Convergence: Tulu, Zephyr\\Recipe \& Distillation};

% Arrows
\draw[arrow] (flan) -- (igpt);
\draw[arrow] (igpt) -- (si);
\draw[brancharrow] (si) -- (brA);
\draw[brancharrow] (si) -- (brB);
\draw[arrow] (si) -- (flanv2);
\draw[arrow] (flanv2) -- (phi1);
\draw[brancharrow] (brA) -- (conv);
\draw[brancharrow] (brB) -- (conv);
\draw[arrow] (phi1) -- (conv);
\end{tikzpicture}
\caption{Evolution path of Thread 1: from the proposal of instruction tuning to the convergence of data engineering methodologies.
Red nodes represent landmark works; gray nodes represent important branches.}
\label{fig:sft-evolution}
\end{figure}

\begin{table}[!ht]
\centering
\caption{Comparison of core methods in Thread 1. The Data Scale column reflects the data scale during the instruction tuning stage.}
\label{tab:sft-comparison}
\small
\begin{tabularx}{\textwidth}{l c >{\raggedright\arraybackslash}p{2.2cm} >{\raggedright\arraybackslash}X >{\raggedright\arraybackslash}X}
\toprule
\textbf{Method} & \textbf{Year} & \textbf{Data Scale} & \textbf{Key Insight} & \textbf{Limitation} \\
\midrule
FLAN & 2021 & 62 tasks, 10 templates each & Zero-shot generalization scales with task diversity & Human-crafted templates; NLP benchmarks only \\
InstructGPT & 2022 & 13K demos & 1.3B SFT+RLHF $>$ 175B GPT-3 in human eval & 40 professional labelers; high cost \\
Self-Instruct & 2022 & 52K from 175 seeds & LM can bootstrap its own instruction data & Style mimicry without substance \\
LIMA & 2023 & 1K curated & Quality $\gg$ quantity (Superficial Alignment Hypothesis) & Single-turn QA only; 65B model \\
Flan-v2 & 2022 & 1,836 tasks & 0.2\% pre-training compute; cross-architecture & Marginal returns after $\sim$282 tasks \\
phi-1 & 2023 & $<$7B tokens (textbook quality) & Data quality can break scaling laws & Code domain only \\
Zephyr & 2023 & GPT-4 distilled & 7B distilled surpasses 70B open models & Dependent on proprietary teacher \\
PEAR & 2026 & Importance-weighted SFT & Best SFT $\neq$ best RL initialization & Early empirical finding \\
\bottomrule
\end{tabularx}
\end{table}

The evolution of Thread 1 follows a clear trajectory:
\begin{enumerate}[leftmargin=1.5em]
  \item \textbf{Problem formulation} (2021): FLAN and T0 demonstrated that multi-task instruction tuning
    can unlock the zero-shot generalization capability of pre-trained LMs.
  \item \textbf{Paradigm establishment} (2022): InstructGPT embedded SFT within a broader alignment pipeline,
    establishing the ``SFT $\rightarrow$ RM $\rightarrow$ PPO'' three-stage training paradigm.
  \item \textbf{Path divergence} (2022--2023): Driven by the problem of data acquisition cost,
    the field diverged into two paths: synthetic data (Self-Instruct $\rightarrow$ Alpaca $\rightarrow$ WizardLM)
    and data curation (LIMA $\rightarrow$ DEITA $\rightarrow$ IFD).
  \item \textbf{Scaling and specialization} (2022--2023): Flan-v2 pushed multi-task scaling to its extreme,
    phi-1 demonstrated that data quality can break scaling laws,
    and domain-specific SFT expanded across code, math, multilingual, and other domains.
  \item \textbf{Convergence and fusion} (2023--): Works such as Tulu and Zephyr began to systematically integrate
    the experience from multiple paths, seeking reproducible optimal training recipes.
\end{enumerate}

\noindent
The core takeaway of this chapter is:
\textbf{SFT is fundamentally a data engineering problem}.
Whether it is Self-Instruct's synthetic data pipeline,
LIMA's Superficial Alignment Hypothesis,
phi-1's textbook-quality data filtering,
or Flan-v2's large-scale task mixing---all
breakthroughs revolve around the central question of how to construct optimal training data.
The model architecture and training algorithm (standard cross-entropy loss on response tokens)
have remained virtually unchanged throughout this evolution;
what has changed is always the \textbf{source}, \textbf{quality}, \textbf{diversity}, and \textbf{scale} of the data.

This insight also points to deep connections between SFT and other threads:
Thread 2 (preference optimization) supplements SFT's binary supervision by introducing preference signals;
Thread 4 (reasoning) extends SFT's training objectives through CoT and process rewards;
Thread 7 (efficiency) makes SFT feasible in resource-constrained scenarios through methods such as LoRA.
These cross-thread connections will be elaborated in subsequent chapters.

\paragraph{Preview of the next chapter.}
SFT trains models using binary supervision (correct vs.\ incorrect),
but human preferences inherently form a continuous spectrum---the
distinction between ``a good answer'' and ``a better answer'' is far beyond what SFT's cross-entropy loss can express.
Thread~2 will introduce preference optimization,
tracing the evolutionary path from comparison-based reward learning through RLHF and DPO to RLVR,
demonstrating how post-training has progressed from ``imitation learning'' to ``preference alignment''.
